\documentclass[a4paper,11pt]{article}
\usepackage[utf8]{inputenc}
\usepackage[T1]{fontenc}
\usepackage[french]{babel}
\usepackage[left=1cm,right=1cm,top=.5cm,bottom=0cm]{geometry}
\usepackage{enumitem}
\usepackage{hyperref}
\usepackage{xcolor}
\usepackage{titlesec}
\usepackage{fontawesome5}
\usepackage{lmodern}
\usepackage[normalem]{ulem}

% --- Couleurs Sanofi ---
\definecolor{primary}{RGB}{102, 45, 145}
\definecolor{secondary}{RGB}{80, 80, 80}

\hypersetup{colorlinks=true, linkcolor=primary, filecolor=primary, urlcolor=primary}

\titleformat{\section}{\large\bfseries\color{primary}\uppercase}{}{0em}{}[\titlerule]
\titlespacing{\section}{0pt}{8pt}{5pt}

\newcommand{\resumeItem}[1]{\item\small{#1}}

\begin{document}

% --- EN-TÊTE ---
\begin{center}
    {\Large \textbf{Loïc Aron MBASSI EWOLO}} \\ \vspace{4pt}
    \textbf{\large Alternance -- Data Analyse / Machine Learning} \\
    \textit{\small Disponible dès septembre 2026} \\ \vspace{4pt}
    \small
    \faMapMarker\ Cergy, France (Mobile IDF) \quad
    \faPhone\ (+33) 7 59 52 33 98 \quad
    \faEnvelope\ \href{mailto:loicaron.mbassiewolo@ensea.fr}{loicaron.mbassiewolo@ensea.fr} \\ \vspace{2pt}
    \faLinkedin\ \href{https://www.linkedin.com/in/loic-mbassi/}{linkedin.com/in/loic-mbassi} \quad
    \faGithub\ \href{https://github.com/Nameless0l}{github.com/Nameless0l} \quad
    \faGlobe\ \href{https://mbassiloic.tech}{mbassiloic.tech}
\end{center}

\vspace{-6pt}

% --- PROFIL ---
\noindent\small{Pipelines Python d'analyse de données médicales à grande échelle, classification d'anomalies cardiaques par ML et tableaux de bord de visualisation en santé : trois axes pratiqués en contexte de recherche (INRIA Saclay / campus IP Paris, UROP avec mentors Google DeepMind, 1\textsuperscript{er} prix IndabaX Africa 2025 -- IA \& Santé). Formé en mathématiques appliquées, génie logiciel et IA (ENSEA / Polytechnique Yaoundé), je suis disponible dès septembre 2026 pour rejoindre les équipes data de Sanofi.}

\vspace{2pt}

% --- FORMATION ---
\section{Formation}
\begin{itemize}[leftmargin=*, label=\tiny$\bullet$, nosep]
    \item \textbf{ENSEA -- École Nationale Supérieure de l'Électronique et de ses Applications} \hfill \textit{2025 -- 2027} \\
    \small{Double diplôme d'Ingénieur -- Systèmes Embarqués, Signal, IA. Cursus : Deep Learning, traitement du signal, probabilités, statistiques.}
    \item \textbf{École Polytechnique de Yaoundé (1\textsuperscript{ère} école d'ingénieur CEMAC)} \hfill \textit{2021 -- 2025} \\
    \small{Diplôme d'Ingénieur de Conception (Master 2) : \textbf{Mathématiques Appliquées}, Génie Logiciel, Algorithmique. Algèbre linéaire (SVD), probabilités, analyse numérique.}
\end{itemize}

% --- COMPÉTENCES ---
\section{Compétences Techniques}
\begin{itemize}[leftmargin=*, nosep]
    \item \small{\textbf{Data \& Machine Learning :} Python, pandas, NumPy, scikit-learn, PyTorch, TensorFlow/Keras, NLP (BERT/Transformers)}
    \item \small{\textbf{Visualisation et reporting :} Streamlit, Flask, matplotlib, seaborn, Power BI (notions)}
    \item \small{\textbf{Collecte et traitement :} SQL, API REST, Git, Docker, nettoyage et normalisation de données, traitement de données textuelles}
    \item \small{\textbf{Langues :} Français (natif), Anglais (scientifique, rédaction et lecture)}
\end{itemize}

% --- EXPÉRIENCES ---
\section{Expériences et Projets}
\begin{itemize}[leftmargin=*, label={-}]

\item \textbf{Stage Recherche -- INRIA Saclay, Équipe TRIBE (campus IP Paris)} \hfill \textit{Avr. -- Août 2026} \\
\textit{Palaiseau -- Programme national PEPR MOBIDEC, collab. École des Ponts ParisTech} \hfill \textit{Python, NLP, pandas}
\vspace{-2pt}
\begin{itemize}[leftmargin=0.4cm, nosep, label=\tiny$\bullet$]
    \item \small{Analyse NLP à grande échelle des politiques de confidentialité d'applications mobiles (BERT/Transformers) : évaluation automatisée de la transparence, de la spécificité et de la conformité RGPD sur plusieurs centaines de documents.}
    \item \small{Conception d'un modèle de score de confiance interprétable, synthétisant plusieurs analyses (NLP, statique, permissions) en un indicateur lisible. Contribution à des publications scientifiques.}
\end{itemize}
\vspace{3pt}

\item \textbf{Assistant de Recherche -- MILA, Institut Québécois d'IA} \hfill \textit{Juin -- Sept. 2025 (Remote)} \\
\textit{Sujet : Grokking et interprétabilité des réseaux de neurones} \hfill \textit{Python, PyTorch, mathématiques}
\vspace{-2pt}
\begin{itemize}[leftmargin=0.4cm, nosep, label=\tiny$\bullet$]
    \item \small{Étude du phénomène de généralisation tardive (grokking) sur des MLP : reproduction d'expériences SOTA, analyse mathématique de la convergence et des dynamiques d'apprentissage.}
    \item \small{Rédaction scientifique en anglais. Collaboration à distance avec des chercheurs du MILA (un des instituts IA de référence au monde).}
\end{itemize}
\vspace{3pt}

\item \textbf{UROP 2024 -- Analyse automatique d'ECG (Santé)} \hfill \textit{2024} \\
\textit{Projet de recherche -- Mentors Google DeepMind} \hfill \textit{TensorFlow/Keras, Computer Vision, traitement du signal}
\vspace{-2pt}
\begin{itemize}[leftmargin=0.4cm, nosep, label=\tiny$\bullet$]
    \item \small{Conception d'un modèle de détection et classification d'anomalies cardiaques à partir d'images d'électrocardiogrammes (ECG sur papier photographié) : extraction de features signal, filtrage, pipeline d'inférence complet.}
    \item \small{Collaboration avec des mentors Google DeepMind. Traitement de données médicales bruyantes et hétérogènes.}
\end{itemize}
\vspace{3pt}

\item \textbf{1\textsuperscript{er} Prix IndabaX Africa 2025 -- IA \& Santé} \hfill \textit{2025} \\
\textit{Hackathon continental IA (données médicales)} \hfill \textit{Python, Flask, Streamlit, API REST}
\vspace{-2pt}
\begin{itemize}[leftmargin=0.4cm, nosep, label=\tiny$\bullet$]
    \item \small{Nettoyage et normalisation d'un jeu de données médicales réelles (valeurs manquantes, anomalies, formats hétérogènes). Entraînement et déploiement d'un modèle ML via API Flask.}
    \item \small{Tableau de bord Streamlit pour la visualisation et l'interprétation des résultats par des non-techniciens. Premier prix continental.}
\end{itemize}
\vspace{3pt}

\item \textbf{Women Safe -- Détection de contenus sexistes (NLP)} \hfill \textit{2023, Compétition MLPC} \\
\textit{Classification de texte -- Réseaux sociaux et offres d'emploi} \hfill \textit{BERT, GPT-2, Transformers}
\vspace{-2pt}
\begin{itemize}[leftmargin=0.4cm, nosep, label=\tiny$\bullet$]
    \item \small{Modèle NLP de détection de misogynie (BERT, GPT-2). Étude des biais éthiques dans les données d'entraînement. Évaluation des faux positifs et implications réglementaires.}
\end{itemize}
\vspace{3pt}

\end{itemize}

% --- DISTINCTIONS ---
\section{Distinctions et Engagement}
\begin{itemize}[leftmargin=*, nosep, label=\tiny$\bullet$]
    \item \small{\textbf{1\textsuperscript{er} prix} IndabaX Africa 2025 (IA \& Santé) \quad \textbf{1\textsuperscript{er} prix} CITS National Hackathon (75 équipes) \quad \textbf{1\textsuperscript{er} prix} Mountain Hub Regional 2024}
    \item \small{Lead Cellule IA, ENSPY : +50\% membres, collaboration Google DeepMind. Animation ateliers ML, NLP, fine-tuning LLM. Articles techniques publiés sur Medium.}
\end{itemize}

\end{document}
